\begin{frame}{확인 문제 (2)}
  \begin{itemize}
    \item \textbf{3.} 같은 데이터를 (a) 무작위 분할, (b)
      환자(그룹) 분할로 나눠 kNN을 학습하면 (a)의 test RMSE가
      비현실적으로 낮아진다. (a)가 왜 좋은 성적을 내는가?
      \vskip0.3em
      {\footnotesize \textbf{답}: (a)에서는 같은 환자의 여러 행이
      train/test에 흩어져, kNN이 test 행의 k-이웃을 \textbf{같은
      환자의 train 행}으로 찾아 ``환자를 인식''한다 --- 관계를
      배운 것이 아니라 환자를 \textbf{외운} 것. 질문: ``그룹 키
      (환자 ID)가 있는가? \texttt{GroupKFold}로 재분할하면?''}
    \item \textbf{4.} ``GBDT $n_{\text{estimators}}$=300을 val
      곡선으로 정했다''는데, val F1이 300에서 \textbf{계속
      상승 중}이다. 한계는?
      \vskip0.3em
      {\footnotesize \textbf{답}: 상승 중이면 ``300이 최적''이
      아니라 ``관찰 범위 내 최선''. val이 $n\approx150$에서 정점
      후 하락했다면 300은 train F1=1.0 포화 상태에서 \textbf{잔차
      (노이즈)를 쫓다 overfit한} 구간 --- 150 부근(early
      stopping)을 골랐어야 함. 기준은 언제나 \textbf{val 정점}.}
  \end{itemize}
\end{frame}
